\section{Geometric collisions at Stieltjes index zero}
\label{gc:section}

The periodic contact reports \cite{RepoPeriodic,RepoContact,RepoMixed} develop
two-factor and spectral completions. Here the shifts themselves move toward
collision. The result covers every finite number of factors and every
argument-derivative order at Stieltjes index zero.

Set $f(x)=\gamma_0(x)=-\psi(x)$ and $h(x)=f(1+x)$.  The Hurwitz shift
identity gives the exact local decomposition
\begin{equation}\label{gc:one-sided}
 f(\{x\})=\frac{\boldsymbol{1}_{x>0}}x+h(x),\qquad -1<x<1,
 \quad x\ne0.
\end{equation}
Here the negative half of the formula uses $\{x\}=1+x$.  The germ $h$ is
holomorphic on $|x|<1$, with
\begin{equation}\label{gc:hseries}
 h(x)=\gamma+\sum_{n\ge1}(-1)^n\zeta(n+1)x^n.
\end{equation}
The finite-part convention is fixed throughout: at each right-hand
singular point $a$ remove $(a,a+\eta)$ and take the constant term in the
\emph{ambient coordinate} $\eta$.  Independent cutoffs are used at distinct
points.  No collision-dependent rescaling of these coordinates is made.

For distinct sufficiently small real shifts $s_1,\ldots,s_d$, define
\begin{equation}\label{gc:Cdef}
 C_d(\boldsymbol{s})=
 \operatorname{FP}\int_{\mathbb R/\mathbb Z}
       \prod_{i=1}^d f(\{x+s_i\})\,dx,
 \qquad
 Q_d=\operatorname{FP}\int_0^1 f(x)^d\,dx.
\end{equation}
For a nonempty subset $S\subseteq\{1,\ldots,d\}$, let $m(S)$ denote its
index with smallest shift.  It is a fixed index when an ordering chamber
is fixed.

\begin{theorem}[Complete local counterterm for any number of factors]
\label{gc:main}
Define
\begin{equation}\label{gc:T}
 T_d(\boldsymbol{s})=
 -\sum_{\substack{S\subseteq\{1,\ldots,d\}\\ |S|\ge2}}
  \ \sum_{i\in S\setminus\{m(S)\}}
  \frac{\displaystyle\prod_{j\notin S}h(s_j-s_i)}
       {\displaystyle\prod_{k\in S\setminus\{i\}}(s_k-s_i)}
       \log(s_i-s_{m(S)}).
\end{equation}
On each fixed ordering chamber,
$C_d(\boldsymbol{s})-T_d(\boldsymbol{s})$ is the restriction of a real
analytic function on a full neighborhood of $\boldsymbol{s}=0$.  Its value
at total collision is $Q_d$.  Consequently
\begin{equation}\label{gc:anypath}
 C_d(\boldsymbol{s})-T_d(\boldsymbol{s})\longrightarrow Q_d
 \qquad (\boldsymbol{s}\longrightarrow0)
\end{equation}
along every path remaining in that chamber, without a lower bound on the
ratios of the gaps.
\end{theorem}

\begin{proof}
Choose $\delta>0$ small enough that all required germs are holomorphic on
$|z|<3\delta$, and restrict the shifts to $|s_i|<\delta/8$.  Use the circle
coordinate $-\delta<x<1-\delta$.  Outside $(-\delta,\delta)$ the entire
integrand is jointly analytic in the shifts, and its integral is analytic.
On $(-\delta,\delta)$, expand the exact product of
\eqref{gc:one-sided} by the subset $S$ of factors contributing their pole.
The empty subset is analytic.  For a nonempty $S$, the product of the
indicator functions is $\boldsymbol{1}_{x>-s_{m(S)}}$, and the contribution
is the finite part of
\begin{equation}\label{gc:subset-integral}
 \int_{-s_{m(S)}}^\delta \frac{q_S(x;\boldsymbol{s})}
                                      {D_S(x;\boldsymbol{s})}\,dx,
 \quad
 q_S(x;\boldsymbol{s})=\prod_{j\notin S}h(x+s_j),
 \quad D_S(x;\boldsymbol{s})=\prod_{i\in S}(x+s_i).
\end{equation}

Let $P_S$ be the interpolation polynomial of degree less than $|S|$ which
agrees with $q_S$ at the nodes $-s_i$, $i\in S$.  Both $P_S$ and
$R_S=(q_S-P_S)/D_S$ extend holomorphically when nodes coalesce.  One
explicit justification is the fixed-contour expression
\begin{equation}\label{gc:hermite}
 P_S(x;\boldsymbol{s})=\frac1{2\pi i}\oint_{|z|=2\delta}
 \frac{q_S(z;\boldsymbol{s})}{D_S(z;\boldsymbol{s})}
 \frac{D_S(z;\boldsymbol{s})-D_S(x;\boldsymbol{s})}{z-x}\,dz.
\end{equation}
Its integrand is holomorphic in the parameters and the quotient of
polynomials is a polynomial in $x$ of the required degree.  Cauchy's
formula gives the jointly holomorphic remainder
\[
 R_S(x;\boldsymbol{s})=\frac1{2\pi i}\oint_{|z|=2\delta}
 \frac{q_S(z;\boldsymbol{s})}
 {D_S(z;\boldsymbol{s})(z-x)}\,dz,
 \qquad |x|<2\delta.
\]

For distinct nodes the rational principal part is
\[
 \frac{P_S(x;\boldsymbol{s})}{D_S(x;\boldsymbol{s})}
 =\sum_{i\in S}\frac{b_{S,i}(\boldsymbol{s})}{x+s_i},
 \qquad
 b_{S,i}(\boldsymbol{s})=
 \frac{q_S(-s_i;\boldsymbol{s})}
      {\prod_{k\in S\setminus\{i\}}(s_k-s_i)}.
\]
Its finite-part integral in \eqref{gc:subset-integral} equals
\begin{equation}\label{gc:integrated}
 \sum_{i\in S}b_{S,i}\log(\delta+s_i)
 -\sum_{i\in S\setminus\{m(S)\}}
       b_{S,i}\log(s_i-s_{m(S)}).
\end{equation}
At the minimal node the discarded primitive is $b_{S,m(S)}\log\eta$,
whose constant term is zero under the specified coordinate convention.
The second sum in \eqref{gc:integrated} is precisely the corresponding
summand of \eqref{gc:T}.  Although the individual coefficients $b_{S,i}$
have poles at coincident nodes, the first sum has the holomorphic form
\begin{equation}\label{gc:upper-contour}
 \frac1{2\pi i}\oint_{|z|=\delta/2}
 \frac{q_S(z;\boldsymbol{s})}{D_S(z;\boldsymbol{s})}
                      \operatorname{Log}(\delta-z)\,dz.
\end{equation}
The logarithm is the holomorphic branch equal to $\log\delta$ at $z=0$.
This contour is deliberately smaller than the interpolation contour in
\eqref{gc:hermite}.  The integral of $R_S$ from $-s_{m(S)}$ to $\delta$ is
also holomorphic, since its lower endpoint is a fixed coordinate function
on the chosen chamber.

At $\boldsymbol{s}=0$, $D_S=x^{|S|}$ and $P_S$ is the Taylor polynomial of
$q_S$ through degree $|S|-1$.  Formula \eqref{gc:upper-contour} is exactly
the finite-part integral of $P_S/x^{|S|}$ from $0$ to $\delta$; the
remainder gives the ordinary integral of the analytic part.  Summing the
subsets therefore gives the same local finite part as the directly
merged product.  The outer integral supplies its remaining part, proving
that the analytic remainder has value $Q_d$.
\end{proof}

\begin{remark}[Meaning of the chamber qualification]
The extension asserted above is an extension \emph{from each chamber}.
Different chambers choose different lower-endpoint indices in
\eqref{gc:subset-integral}; no assertion that the resulting analytic germs
agree off the total diagonal is needed.  All their values at total
collision are the prescribed merged finite part.  In particular the
theorem proves a path-independent limit only after subtracting the full
$T_d$, including any finite constant it contributes on the chosen path.
\end{remark}

\subsection{Every argument-derivative order}
For nonnegative integers $r_1,\ldots,r_d$, put
\[
 d_i=r_i+1,\qquad A_i=(-1)^{r_i}r_i!,\qquad
 h_i(x)=\gamma_0^{(r_i)}(1+x).
\]
Ordinary differentiation away from the integer endpoints gives
\[
 \gamma_0^{(r_i)}(\{x\})
    =A_i\boldsymbol{1}_{x>0}x^{-d_i}+h_i(x).
\]
Define $C_{\boldsymbol r}(\boldsymbol s)$ and $Q_{\boldsymbol r}$ by
replacing the factors in \eqref{gc:Cdef} by the indicated argument
derivatives.  For a nonempty $S$, set
\[
 A_S=\prod_{i\in S}A_i,\qquad
 q_S(x;\boldsymbol s)=\prod_{j\notin S}h_j(x+s_j),\qquad
 D_S(x;\boldsymbol s)=\prod_{i\in S}(x+s_i)^{d_i}.
\]
For $i\in S$ and $1\le q\le d_i$, let
\begin{equation}\label{gc:B}
 B_{S,i,q}(\boldsymbol s)=
 \frac1{(d_i-q)!}
 \left.\left(\frac{d}{dx}\right)^{d_i-q}
 \frac{q_S(x;\boldsymbol s)}
 {\prod_{k\in S\setminus\{i\}}(x+s_k)^{d_k}}
 \right|_{x=-s_i}.
\end{equation}

\begin{theorem}[Repeated-pole collision formula]\label{gc:derivative}
Let
\begin{align}\label{gc:Tr}
 T_{\boldsymbol r}(\boldsymbol s)
 ={}&\sum_{\substack{S\subseteq\{1,\ldots,d\}\\ |S|\ge2}}
 A_S\sum_{i\in S\setminus\{m(S)\}}
 \left\{-B_{S,i,1}\log(s_i-s_{m(S)})\right.\notag\\[-2pt]
 &\hspace{42mm}\left.
 +\sum_{q=2}^{d_i}\frac{B_{S,i,q}}{q-1}
                   (s_i-s_{m(S)})^{1-q}\right\}.
\end{align}
Then $C_{\boldsymbol r}-T_{\boldsymbol r}$ has a real analytic extension
from each ordering chamber through total collision, with value
$Q_{\boldsymbol r}$ there.
\end{theorem}

\begin{proof}
Repeat the preceding proof with Hermite interpolation of multiplicity
$d_i$ at $-s_i$.  Formula \eqref{gc:hermite}, with the new $D_S$, proves
holomorphy of the interpolation polynomial and its remainder, including
coalescing nodes.  The principal part is
$\sum_{i\in S}\sum_{q=1}^{d_i}B_{S,i,q}(x+s_i)^{-q}$.  Its upper primitive
is again \eqref{gc:upper-contour}, multiplied by $A_S$, since
\[
 \frac1{(q-1)!}\frac{d^{q-1}}{dz^{q-1}}
      \operatorname{Log}(\delta-z)
 =-\frac{(\delta-z)^{1-q}}{q-1},\qquad q\ge2.
\]
At a nonminimal node, subtracting the lower primitive gives exactly the
bracket in \eqref{gc:Tr}.  At the minimal node each lower primitive is a
pure power of $\eta$ or a multiple of $\log\eta$, and hence has zero
constant term.  The same coalescence argument identifies the value of the
analytic remainder with $Q_{\boldsymbol r}$.
\end{proof}

This proof differentiates only ordinary local analytic functions in
\eqref{gc:B}.  Differentiating an already extended periodic distribution
can generate contact terms, and does not supply a substitute for the
direct repeated-pole proof.  For example, with two factors and shifts
$(0,a)$, the formula gives
\[
 T_{(1,0)}(0,a)=\frac{\log a}{a^2},\qquad
 T_{(0,1)}(0,a)=\frac{1-\log a}{a^2}.
\]
Both analytic remainders tend to
$Q_{(1,0)}=Q_{(0,1)}=\zeta(2)$, which also follows by integrating
$ff'=(f^2)'/2$ and taking the endpoint finite part.
The separated integrals satisfy the exact identity
\begin{equation}\label{gc:pair-sum}
 C_{(1,0)}(0,a)+C_{(0,1)}(0,a)
   =\frac{\pi^2}{\sin^2(\pi a)},\qquad 0<a<1.
\end{equation}
Indeed, apply the fundamental theorem of calculus to the product
$f(x)f(\{x+a\})$ on the two nonsingular intervals.  At $0$ the left
boundary value minus the constant in the right Laurent expansion is
$-f'(a)$; at $1-a$ it is $-f'(1-a)$.  Taking the finite parts and using
trigamma reflection gives
$-f'(a)-f'(1-a)=\psi'(a)+\psi'(1-a)=\pi^2/\sin^2(\pi a)$.
Thus \eqref{gc:pair-sum} records the boundary correction required by the
specified finite part.  The sum of the two local counterterms is $a^{-2}$,
and the sum of their limiting remainders is $2\zeta(2)=\pi^2/3$, in
agreement with the Taylor expansion of its right side.

\subsection{Fixed rates and a finite zeta--logarithm expression}
For $s_i=a_i\epsilon$ with fixed distinct real rates, define
$m(S)=\operatorname*{argmin}_{i\in S}a_i$.  Coefficient extraction from
\eqref{gc:T} gives
\begin{align}\label{gc:fixed-constant}
 \operatorname{CT}_{\epsilon\downarrow0}C_d(\epsilon\boldsymbol a)
 ={}&Q_d-\sum_{\substack{S\subseteq\{1,\ldots,d\}\\ |S|\ge2}}
 \ \sum_{i\in S\setminus\{m(S)\}}
 \frac{\log(a_i-a_{m(S)})}
      {\prod_{k\in S\setminus\{i\}}(a_k-a_i)}\notag\\
 &\hspace{18mm}\cdot
 [\epsilon^{|S|-1}]
       \prod_{j\notin S}h\bigl(\epsilon(a_j-a_i)\bigr).
\end{align}
Here the constant-term operation discards positive powers of
$\log\epsilon$ even when their power of $\epsilon$ is zero.  Only finitely
many Taylor coefficients of \eqref{gc:hseries} occur.  In particular the
finite anomaly is an explicit polynomial in Euler's constant and ordinary
zeta values, with rational functions of the rates and logarithms of their
positive differences as coefficients.  This is a local conclusion; it
does not evaluate the merged global constant $Q_d$.

For equally spaced triples, \eqref{gc:fixed-constant} recovers the earlier
three-factor contact specialization \cite{RepoMixed}:
\[
 \operatorname{CT}_{\epsilon\downarrow0}
 C_3(0,\epsilon,2\epsilon)=Q_3+\frac{\zeta(2)}2\log2.
\]
For four equally spaced points, the new formula is
\begin{equation}\label{gc:quartic-example}
 \operatorname{CT}_{\epsilon\downarrow0}
 C_4(0,\epsilon,2\epsilon,3\epsilon)
 =Q_4+\gamma\zeta(2)(2\log2+\log3)
       -\frac32\zeta(3)(3\log2+\log3).
\end{equation}

\subsection{A hierarchical collision}
Take shifts $(0,\epsilon,\epsilon+\epsilon^2)$, so the last pair approaches
one another on a smaller scale than its distance from the first point.
The four terms in \eqref{gc:T} are exactly
\begin{align*}
 T_{\{1,2,3\}}&=\frac{\log\epsilon}{\epsilon^3}
 -\frac{\log(\epsilon+\epsilon^2)}{\epsilon^3(1+\epsilon)},\\
 T_{\{1,2\}}&=\frac{h(\epsilon^2)}\epsilon\log\epsilon,\\
 T_{\{1,3\}}&=\frac{h(-\epsilon^2)}{\epsilon+\epsilon^2}
                         \log(\epsilon+\epsilon^2),\\
 T_{\{2,3\}}&=\frac{h(-\epsilon-\epsilon^2)}{\epsilon^2}
                         \log(\epsilon^2).
\end{align*}
Writing $\Lambda=\log\epsilon$, expanding these exact terms, and applying
Theorem~\ref{gc:main} yields
\begin{align}\label{gc:hierarchy}
 C_3(0,\epsilon,\epsilon+\epsilon^2)
 ={}&\frac{(1+2\gamma)\Lambda-1}{\epsilon^2}
 +\frac{(2\gamma+2\zeta(2)-1)\Lambda+\tfrac32}{\epsilon}\notag\\
 &+(1-\gamma+2\zeta(2)+2\zeta(3))\Lambda
 +Q_3+\gamma-\frac{11}{6}
 +O(\epsilon|\log\epsilon|).
\end{align}
The finite constant is therefore $Q_3+\gamma-11/6$.  The rational term
$-11/6$ comes from the all-pole subset, while the term $\gamma$ comes from
$T_{\{1,3\}}$.  This example shows why one cannot obtain hierarchical
constant terms by retaining only the singular powers of a fixed-rate
expansion and subsequently varying its rate parameters.

\subsection{Verification and precise scope}
The companion script \texttt{verify\_collision.py} verifies the quartic
constant and the full hierarchical expansion through order zero by exact
symbolic algebra.  It also evaluates the separated circle integral by
subtracting its actual right-hand poles interval by interval, independently
of \eqref{gc:T}.  Direct Laurent subtraction at
the merged endpoint gives
\begin{align*}
 Q_3&=-1.9662627343915343406643753384116\ldots,\\
 Q_4&=\phantom{-}13.510120376015827411606240712830\ldots.
\end{align*}
The computed differences $C_d-T_d-Q_d$ tend to zero linearly on both
displayed paths, at four values of $\epsilon$ for each path and $55$
working decimal digits.  The separate script
\texttt{verify\_derivative\_pair.py} checks both displayed derivative
counterterms at the same four gaps, using independent double-pole
subtraction and $65$ working decimal digits.  Its four comparisons with
\eqref{gc:pair-sum} have observed absolute residual at most
$6.4\times10^{-51}$.  The package includes the raw JSON records.
These floating-point diagnostics check implementation and signs; the
interpolation proof supplies the theorem.

The result settles the arbitrary-factor and arbitrary argument-derivative
parts of the unit-frequency geometric collision question at Stieltjes index zero,
including hierarchical approaches.  Higher Stieltjes indices introduce
logarithmic singular germs and are not covered by the rational
interpolation argument.  Associativity of successively renormalized
partial mergers also remains a separate question: an exact simultaneous
subtraction does not by itself identify every iterated subtraction rule.
Two further extensions are a complete statement for integer-frequency
singular grids, with each local pole amplitude scaled in the ambient
coordinate, and a weighted version for analytic periodic test functions
which would provide collision formulas for Fourier moments.  Neither is
included in the theorems proved here.
